%========================================================
% CHAPTER 3 — Numerical Experiments and Results
%========================================================

The theoretical framework developed in the previous chapters is implemented and
investigated numerically in this chapter. The objective is to move from the
analytical description of polarization-entangled photon propagation to a
computational framework capable of quantifying how physically relevant fiber and
source parameters affect the distributed quantum state.

The numerical analysis is organized around a progressive increase in physical
complexity. Controlled limiting cases are considered first in order to validate the
implementation against analytical predictions. The study is then extended to
frequency-dependent birefringent propagation, two-arm polarization-mode
dispersion, and statistical ensembles of randomly birefringent fibers. Particular
attention is given to the dependence of the output state on quantities that can be
related directly to the physical system, including photon bandwidth, differential
group delay, fiber length, birefringence statistics, principal-state orientations,
and the spectral correlations of the photon pair.

The purpose of the numerical experiments is therefore not limited to verifying the
correctness of the implementation. The simulator is used as an analysis tool to
identify the physical regimes in which fiber propagation produces negligible,
partial, or substantial degradation of polarization entanglement, and to determine
how these regimes are reflected in experimentally accessible quantities. The
output states are characterized through polarization correlations, reduced density
operators, purity, fidelity, concurrence, and Bell-nonlocality indicators. When the
fiber model is stochastic, these quantities are treated statistically over an
ensemble of fiber realizations rather than solely through their mean values.

A further objective is to establish a direct connection between the numerical model
and the experimental investigation developed later in this work. Whenever
possible, the numerical experiments are therefore formulated in terms of parameters
that can either be controlled or estimated experimentally. This allows the
simulator to provide reference predictions for the interpretation of measured
single-photon and coincidence statistics and, more generally, to investigate the
range of fiber and source conditions compatible with the preservation of useful
polarization entanglement.

%--------------------------------------------------------
% MATLAB computational framework
%--------------------------------------------------------

\section{Computational Framework and MATLAB Library}
\label{sec:computational_framework_matlab}

All numerical models presented in this chapter are implemented in MATLAB. This
choice allows the simulator to be integrated directly into the pre-existing
computational codebase used throughout the project, avoiding the introduction of a
separate software environment and facilitating maintenance, validation, and future
extensions of the framework.

MATLAB is also particularly well suited to the mathematical structure of the
problem considered in this work. Quantum states, density operators, Jones
operators, Pauli observables, tensor products, and frequency-dependent propagation
matrices are represented directly through complex vectors and matrices. As a
consequence, the numerical objects used in the implementation closely reproduce
both the notation and the algebraic structure of their theoretical counterparts.
For example, single-photon polarization operators are represented by
\(2\times2\) matrices, bipartite states and operators by vectors and matrices in
\(\mathbb{C}^{2}\otimes\mathbb{C}^{2}\), and local bipartite transformations are
constructed explicitly through tensor products.

The resulting implementation is organized as a modular MATLAB library rather than
as a collection of experiment-specific scripts. Reusable routines are provided for
state preparation, operator construction, local and frequency-dependent
propagation, quantum-channel application, partial traces, correlation evaluation,
state metrics, spectral integration, segmented-fiber construction, statistical
sampling, and numerical validation. Higher-level numerical experiments then
combine these validated components without redefining the underlying mathematical
operations.

This architecture provides a direct correspondence between the theoretical model
and its computational realization:

\[
\begin{array}{c}
\text{Physical and Mathematical Model}
\\
\downarrow
\\
\text{Reusable MATLAB Operators and State Routines}
\\
\downarrow
\\
\text{Frequency-Resolved Fiber Propagation}
\\
\downarrow
\\
\text{Spectral and Statistical Averaging}
\\
\downarrow
\\
\text{Density Operator and Observable Evaluation}
\\
\downarrow
\\
\text{Numerical Experiments and Experimental Predictions}.
\end{array}
\]

Separating the reusable numerical library from the individual experiment
workflows has two main advantages. First, the same implementation of each
mathematical operation is used throughout the thesis, reducing duplication and
making numerical consistency easier to verify. Second, progressively more complex
fiber models can be constructed by composition of previously validated components,
following the same hierarchical strategy adopted in the theoretical development.

The numerical study presented in the following sections therefore proceeds from
controlled validation cases toward increasingly realistic and statistically
characterized models of fiber propagation.

%--------------------------------------------------------
% Chapter sections
%--------------------------------------------------------
% \section{Numerical Methodology Validation}
% \label{sec:03_1_numerical_methodology_validation}
% \input{chapters/numerical_experiments/03_1_numerical_methodology_validation}


